\nuevaReceta{Remojón Tradicional}{45 min}{receta:remojon}
    
    \begin{minipage}[t]{0.35\textwidth}
        \section*{Ingredientes}
        \begin{itemize}[leftmargin=*]
            \item Patatas (aprox. 3-4 medianas)
            \item 1 Cebolla
            \item 1 Tomate maduro
            \item Pimientos secos
            \item Longaniza de calidad
            \item 2 Dientes de ajo
            \item 1 Huevo por persona
            \item Aceite de oliva y sal
        \end{itemize}
    \end{minipage}
    \hfill
    \begin{minipage}[t]{0.6\textwidth}
        \section*{Preparación}
        \begin{enumerate}
            \item \textbf{Cocción base:} Pelar y cortar las patatas. Ponerlas a cocer en una olla con agua, sal, un chorro de aceite, la cebolla entera y el tomate.
            \item \textbf{Pimientos y carne:} Abrir los pimientos secos, quitar las semillas e incorporarlos a la olla junto con la longaniza.
            \item \textbf{Huevos:} En un cazo aparte, cocer los huevos (10 min) y reservar.
            \item \textbf{El majado:} Una vez cocidos, sacar de la olla la cebolla y el tomate. Pasarlos por la batidora junto con los dientes de ajo crudos hasta obtener una mezcla fina.
            \item \textbf{Tratamiento de la carne:} Sacar la longaniza cuando esté cocida y desmenuzarla en trozos pequeños. 
            \item \textbf{Limpieza:} Escurrir casi toda el agua de la cocción (reservarla en un bol). Sacar los pimientos y quitarles la piel cuidadosamente.
            \item \textbf{Unión:} Juntar en la olla las patatas con el batido de cebolla/tomate, la carne de los pimientos ya limpios y la longaniza. Remover bien. Si queda muy seco, añadir un poco del agua de cocción reservada.
            \item \textbf{Toque final:} Poner al fuego suave. Rallar los huevos cocidos por encima y remover mientras termina de ligarse todo en caliente.
        \end{enumerate}
    \end{minipage}

    \vspace{1em}
    \begin{tcolorbox}[colback=orange!5, colframe=orange!50, title=Nota del Chef, size=small]
        El secreto está en el agua de la cocción; no la tires toda de golpe, ya que es la que te dará el punto exacto de cremosidad al final.
    \end{tcolorbox}
\end{receta}
